\section{The Recursive Spawning Protocol}\label{sec:rs-protocol}

\subsection{Overview and Design Objectives}

The Recursive Spawning Protocol (RSP) is the operational mechanism by which the BX3 Framework converts multi-agent spawning from a forensic reconstruction problem into a runtime-verified structural property. Every spawn event is governed by a mandatory five-stage sequence executed across all three BX3 layers. No stage is advisory; no layer may be bypassed; no operational urgency constitutes an exception to any protocol requirement. The result is that every agent in any spawn chain is, by architectural construction, traceable to a human principal, constrained to a scope narrower than its parent, and recorded in an immutable forensic ledger before it enters operational state.

The protocol serves three concurrent objectives. First, it ensures that every child agent carries an immutable parent pointer binding it to its originator for its entire operational lifetime. Second, it enforces that every spawn event passes through three independent validation gates—Purpose Layer, Bounds Engine, and Fact Layer—before the child is activated. Third, it guarantees that the complete chain of any active agent terminates at a human principal within a bounded number of steps, making the chain always traversable by any auditor with ledger access.

\subsection{Five-Stage Spawn Sequence}

Every spawn event under the RSP executes in five mandatory stages, each corresponding to a distinct architectural responsibility. The stages are sequential and atomic: a child agent is not activated until all five stages have completed successfully. Failure at any stage halts the spawn and triggers mandatory refusal logging.

\textbf{Stage 1 — Purpose Layer Validation.} The parent agent submits a spawn request to the Purpose Layer containing the proposed child's purpose declaration: the intended task scope, output type, tool access requirements, and the parent's active charter reference. The Purpose Layer evaluates whether the proposed scope falls within the parent agent's registered purpose charter as established at the parent's activation. If the proposed scope is outside the charter—meaning the child would be authorized to perform actions the parent itself was not chartered to perform—the Purpose Layer rejects the spawn with a \texttt{PURPOSE\_VIOLATION} refusal code. If the scope is within charter, the Purpose Layer issues a signed \texttt{PURPOSE\_ATTESTATION} token: a cryptographic attestation binding the proposed child to the specific purpose clause that justifies its existence. This token is forwarded with the spawn request and is immutable—it cannot be forged, replayed, or altered without detection.

\textbf{Stage 2 — Bounds Engine Depth Check.} The spawn request, now bearing the Purpose Layer attestation, is submitted to the Bounds Engine. The Bounds Engine evaluates two constraints before issuing any approval. The first is the \textbf{max spawn depth} parameter, denoted $D_{\max}$, which is a system-wide configurable constant defining the maximum allowable depth for any agent in the spawn chain. The second is the \textbf{convergence criteria}, which evaluates whether the proposed child represents a genuine task decomposition from its parent rather than lateral replication or depth gaming. If the parent agent's current depth $d$ satisfies $d \geq D_{\max}$, the Bounds Engine refuses the spawn and returns \texttt{DEPTH\_LIMIT\_REACHED}. If the convergence criteria fail—meaning the child's scope is not a proper narrowing of the parent's scope—the Bounds Engine logs a \texttt{NON\_CONVERGENT\_SPAWN} warning and may refuse the spawn if convergence failure indicates scope equivalence. If both constraints pass, the Bounds Engine records the approved spawn depth $d+1$ and forwards the validated request to the Fact Layer.

\textbf{Stage 3 — Fact Layer Immutable Pointer Creation.} Upon receiving the Bounds Engine approval, the Fact Layer generates the child's immutable parent pointer. This pointer is the central structural primitive of the RSP. It is not a mutable attribute—it is a cryptographic, content-addressed identifier derived from the SHA-256 hash of a canonical tuple: $(\text{parent\_id}, \text{purpose\_attestation}, \text{spawn\_timestamp}, \text{child\_public\_key})$. The hash is computed over this tuple at the moment of spawn; the result becomes the child's permanent identity attribute and is never recomputed. This construction ensures that the parent pointer cannot be altered retroactively: any change to the tuple's constituents produces a different hash, breaking the chain's integrity verification. The Fact Layer writes the immutable pointer to the forensic ledger as part of the spawn event record.

\textbf{Stage 4 — Forensic Ledger Recording.} The Fact Layer appends a complete, signed record of the spawn event to the append-only forensic ledger before returning any spawn confirmation. The ledger entry is the authoritative, tamper-evident record of every spawn event in the system. It includes: the child agent's new identifier, the immutable parent pointer (SHA-256 hash), the Purpose Layer attestation token, the approved spawn depth, the convergence evaluation result, a Unix epoch timestamp at millisecond precision, and the identifier of the Fact Layer instance that processed the event. The ledger is append-only: no entry may be deleted, modified, or redacted after inscription. The ledger entry identifier is returned to the parent agent as the child's forensic birth certificate.

\textbf{Stage 5 — Chain Termination Verification and Child Activation.} Before the child agent receives any tool access, state mutation authority, or task delegation, the Fact Layer performs a synchronous chain traversal. Starting from the newly created child, it follows each immutable parent pointer upward—parent to grandparent, and so on—until it reaches a node that does not appear as a child in any ledger entry. That node is the human principal at the root of the chain. The traversal must succeed within $D_{\max}$ steps. If the chain terminates at a human principal within the depth bound, the child is activated with its immutable parent pointer, its approved spawn depth, its constrained purpose scope, and its full lineage chain. If the chain traversal fails to reach a human within the depth bound, the spawn is invalidated, the child is never activated, and a \texttt{CHAIN\_TERMINATION\_FAILED} refusal is logged.

\subsection{Depth Management}

Depth management in the RSP serves two purposes: preventing unbounded resource consumption and ensuring that chain traversal by auditors is bounded by a known finite limit. Depth is managed through the interaction of the Bounds Engine's $D_{\max}$ parameter, the spawn refusal mechanism, and the convergence criteria.

\textbf{Maximum spawn depth.} The system-wide maximum spawn depth, $D_{\max}$, is the absolute upper bound on the depth of any agent in any spawn chain. At initialization, the root human-to-agent spawn event establishes depth $d=0$ for the root agent. Every subsequent agent-to-agent spawn increments the depth counter by one. The Bounds Engine maintains an accurate depth counter for every active agent; this counter is recalculated at each spawn event by querying the forensic ledger for the agent's chain depth rather than relying on any mutable in-memory value. The $D_{\max}$ parameter is enforced as a hard ceiling: no spawn event may result in a child at depth $d \geq D_{\max}$.

\textbf{Spawn refusal at limit.} When a spawn request arrives at the Bounds Engine and the parent agent's current depth $d$ satisfies $d \geq D_{\max}$, the spawn is refused. The parent agent receives a \texttt{DEPTH\_LIMIT\_REACHED} refusal code. No child identity is issued, no tool grants are extended, no task delegation is made, and no ledger entry for a successful spawn is created. The refusal event itself is logged to the forensic ledger with full contextual information. The parent agent may not retry the same spawn without first reducing the active chain depth below $D_{\max}$ through child termination. There is no administrative override, no emergency bypass, and no operational mode in which the depth limit may be exceeded.

\textbf{Convergence criteria.} The convergence criteria prevent depth gaming: the practice of spawning children that perform work equivalent to the parent's own scope purely to inflate chain depth without achieving genuine task decomposition. The criteria evaluate two axes. \textit{Scope narrowing} requires that the child's purpose scope is a proper subset of the parent's active scope—a child performing soil moisture calibration is within scope when the parent performs irrigation planning, but a child performing identical irrigation planning is not. \textit{Task decomposition} requires that the child's task introduces structural distinction from the parent: a different information domain, different tool set, different epistemic object, or different output type. The convergence check is implemented as a deterministic rule-based classifier that evaluates the child's purpose declaration against the parent's active scope descriptor. It produces a binary \texttt{CONVERGENT} or \texttt{NON\_CONVERGENT} decision with an associated confidence score in $[0.0, 1.0]$. Spawns with confidence scores below $0.6$ trigger a \texttt{NON\_CONVERGENT\_SPAWN} escalation and are logged to the forensic ledger for policy review; they may be blocked if the scope analysis indicates equivalence rather than decomposition.

\subsection{Immutable Parent Pointer}

The immutable parent pointer $\alpha(c) = p$ is the foundational structural primitive of the RSP. Its immutability is not a policy convention—it is a structural property enforced by the Fact Layer's append-only write protocol at the moment of spawn. We specify its formal properties.

For any child agent $c$ created via a valid RSP spawn event, the relation $\alpha(c) = p$ is established exactly once and satisfies four governing properties:

\begin{enumerate}
    \item[\textbf{P1 — Uniqueness.}] There exists exactly one $p$ such that $\alpha(c) = p$. No child agent has more than one parent pointer.
    \item[\textbf{P2 — Immutability.}] After the initial Fact Layer write, $\alpha(c)$ cannot be changed by any actor—including the parent $p$, the child $c$, any system administrator, or any other agent. The pointer is encoded in the child's foundational identity and is inseparable from it.
    \item[\textbf{P3 — Verifiability.}] For any agent $c$, the value $\alpha(c)$ can be independently verified by any auditor with forensic ledger access by querying the ledger entry for $c$'s spawn event and recomputing the SHA-256 hash of the canonical tuple.
    \item[\textbf{P4 — Non-repudiation.}] The parent pointer is established atomically with the child's activation. The child does not exist as an accountable actor until the parent pointer write is confirmed by the Fact Layer. If provisioning fails, no parent pointer is written and the child does not exist.
\end{enumerate}

The immutable parent pointer chain is the structure that makes drift structurally impossible. An agent is orphaned only if its parent pointer is severed—which is impossible under P2. An agent is drifted only if its chain does not reach a human—which the Fact Layer's Stage 5 chain verification prevents by construction. An agent operates only if all three layer gates passed—which is logged and verifiable. The conjunction of the three failure modes that constitute the drift triad is thereby made architecturally impossible.

\subsection{Forensic Ledger}

The forensic ledger $\mathcal{F}$ is the append-only, tamper-evident record of all RSP events maintained by the Fact Layer. Every spawn approval and every spawn refusal produces a signed ledger entry before any state change takes effect. The ledger is the authoritative source for all accountability queries: chain reconstruction, termination verification, audit review, and incident investigation.

Each ledger entry is a signed JSON document containing the complete context of the spawn event. The entry schema captures: the entry identifier, event type (spawn approved or spawn refused), child and parent agent identifiers, the immutable parent pointer hash, the Purpose Layer attestation token, the approved spawn depth, the convergence evaluation result, the millisecond-precision timestamp, the Fact Layer instance identifier, the chain termination verification result, and—if applicable—the refusal code and refusal context. Entries are hash-chained: each entry includes the hash of the previous entry, forming a Merkle-chain that makes any historical tampering detectable.

Agents may append entries to the ledger but may not read, modify, or delete existing entries. Read access is reserved for authorized auditors and human principals. The ledger supports three canonical query types: reconstructing the full parent chain for any agent by iterating parent pointers backward to the human anchor; enumerating all spawn events within a given time window; and counting refusal events by type or by agent to identify anomalous spawning behavior.

\subsection{Refusal Handling}

Every refusal condition in the RSP is a first-class protocol event with mandatory forensic logging. Refusal codes are deterministic and exhaustive: no spawn event produces an untyped or ambiguous outcome. The explicit refusal conditions are:

\begin{itemize}
    \item \texttt{PURPOSE\_VIOLATION}: The proposed child's scope falls outside the parent agent's purpose charter. The spawn is refused at Stage 1.
    \item \texttt{DEPTH\_LIMIT\_REACHED}: The parent agent is at or beyond $D_{\max}$. The spawn is refused at Stage 2.
    \item \texttt{NON\_CONVERGENT\_SPAWN}: The convergence criteria indicate that the proposed child's scope is equivalent to or broader than the parent's, constituting depth gaming. The spawn may be refused at Stage 2 depending on configuration.
    \item \texttt{CHAIN\_TERMINATION\_FAILED}: The Fact Layer's chain traversal (Stage 5) failed to locate a human principal within $D_{\max}$ steps. The spawn is invalidated and the child is not activated.
    \item \texttt{FACT\_LAYER\_UNAVAILABLE}: The Fact Layer could not be reached or could not confirm the ledger write. The spawn is refused and logged as a system infrastructure failure.
\end{itemize}

There are no provisional activations, no graceful degradations, and no fallback spawn paths. An agent that receives a refusal code must log the refusal to its own forensic ledger entry and may not re-attempt the same spawn without first resolving the underlying condition. This property—combined with the mandatory refusal logging—ensures that the forensic ledger contains a complete record of all spawn attempts, successful or not, making the system fully auditable even in failure scenarios.

\subsection{Protocol Invariants}

The RSP maintains four structural invariants throughout its execution. These are not runtime conventions; they are properties enforced by the protocol's state machine transitions and verified by the Fact Layer pre-commit hook at every spawn event.

\begin{enumerate}
    \item[\textbf{I1 — No orphan activation.}] Every child agent is activated only after all three layer gates (Purpose Layer, Bounds Engine, Fact Layer) have issued approval. There is no activation path that bypasses any gate.
    \item[\textbf{I2 — Immutable parentage.}] Once a child agent is created, its parent pointer $\alpha(c)$ cannot be altered by any actor for any reason. The parentage relation is permanent.
    \item[\textbf{I3 — Verified chain termination.}] The Fact Layer's chain traversal is synchronous and must succeed before spawn confirmation is issued. Every active agent in the system has a verified chain to a human principal.
    \item[\textbf{I4 — Complete ledger coverage.}] Every spawn event and every refusal event produces a signed entry in the forensic ledger. No event is unlogged.
\end{enumerate}

These four invariants together encode the RSP's accountability guarantee: every agent in every spawn chain is authorized by a human principal, bounded by a finite depth, constrained to a narrower scope than its parent, and recorded in an immutable ledger. The drift triad—orphaned, drifted, and operating simultaneously—is not merely statistically unlikely under these conditions. It is structurally impossible. An agent cannot be orphaned from its parent without violating I2. It cannot be drifted without violating I1 or I3. It cannot be operating without having passed all three layer validations, which are recorded and verifiable under I4. The RSP makes these conditions nonexistent by construction.
